\documentclass[11pt]{article}
\usepackage[italian]{babel}
\newcommand{\imgdir}{images/it}
% Shared layout for the Illuminati manual. Each manual-<lang>.tex loads this
% after \documentclass and sets \manuallang (babel name) and \imgdir.
% Build with XeLaTeX (see build.ps1); the fonts are the app's own.

\usepackage{fontspec}
\setmainfont{Literata}[
  Path=../../assets/fonts/literata/, Extension=.ttf,
  UprightFont=*-Regular, BoldFont=*-Bold, ItalicFont=*-Italic, BoldItalicFont=*-BoldItalic]
\setsansfont{AtkinsonHyperlegible}[
  Path=../../assets/fonts/atkinson/, Extension=.ttf,
  UprightFont=*-Regular, BoldFont=*-Bold, ItalicFont=*-Italic, BoldItalicFont=*-BoldItalic]
\newfontfamily\dyslexicfont{OpenDyslexic}[
  Path=../../assets/fonts/opendyslexic/, Extension=.otf,
  UprightFont=*-Regular, BoldFont=*-Bold, ItalicFont=*-Italic, BoldItalicFont=*-Bold-Italic]
\newfontfamily\garamondfont{EBGaramond}[
  Path=../../assets/fonts/ebgaramond/, Extension=.ttf,
  UprightFont=*-Regular, BoldFont=*-Bold, ItalicFont=*-Italic, BoldItalicFont=*-BoldItalic]

\usepackage[a4paper, margin=2.4cm, headheight=14pt]{geometry}
\usepackage{graphicx}
\usepackage[table]{xcolor}
\usepackage{booktabs, tabularx, array}
\usepackage{enumitem}
\usepackage{fancyhdr}
\usepackage{titlesec}
\usepackage[most]{tcolorbox}
\usepackage[hidelinks]{hyperref}

% Colours from the app's logo.
\definecolor{violet}{HTML}{34206A}
\definecolor{violetlight}{HTML}{EDE8F7}
\definecolor{gold}{HTML}{C9A227}
\definecolor{goldlight}{HTML}{FBF4DF}

\setlength{\parskip}{0.55em}
\setlength{\parindent}{0pt}
\setlist{itemsep=0.25em, topsep=0.3em}
\linespread{1.08}

\titleformat{\section}{\sffamily\Large\bfseries\color{violet}}{\thesection}{0.8em}{}
\titleformat{\subsection}{\sffamily\large\bfseries\color{violet}}{\thesubsection}{0.7em}{}
\titlespacing*{\section}{0pt}{1.6em}{0.6em}
\titlespacing*{\subsection}{0pt}{1.1em}{0.35em}
% Don't leave a heading alone at the bottom of a page.
\usepackage{needspace, etoolbox}
\pretocmd{\section}{\needspace{8\baselineskip}}{}{}
\pretocmd{\subsection}{\needspace{5\baselineskip}}{}{}

\pagestyle{fancy}
\fancyhf{}
\fancyhead[L]{\sffamily\small\color{violet} Illuminati}
\fancyhead[R]{\sffamily\small\color{gray}\nouppercase{\rightmark}}
\fancyfoot[C]{\sffamily\small\thepage}
\renewcommand{\headrulewidth}{0.3pt}
% \rightmark is the first section mark on the page.
\renewcommand{\sectionmark}[1]{\markboth{#1}{#1}}
\renewcommand{\subsectionmark}[1]{}

% Keep catalogue names whole in every language.
\babelhyphenation{Wiki-source Guten-berg}

% A key or key combination, e.g. \kbd{Ctrl+F}.
\newcommand{\kbd}[1]{\tikz[baseline=(k.base)]\node[draw=gray!60, fill=gray!8,
  rounded corners=2pt, inner sep=2pt, font=\sffamily\small](k){#1};}
\usepackage{tikz}

% Arrow keys (the sans font has no arrow characters).
\newcommand{\kright}{\kbd{\ensuremath{\rightarrow}}}
\newcommand{\kleft}{\kbd{\ensuremath{\leftarrow}}}
\newcommand{\kup}{\kbd{\ensuremath{\uparrow}}}
\newcommand{\kdown}{\kbd{\ensuremath{\downarrow}}}

% A button, menu item or label as it appears in the app.
\newcommand{\ui}[1]{\textsf{\textbf{#1}}}

% Screenshot with caption; skipped if the image file isn't there.
\newcommand{\shot}[2]{%
  \IfFileExists{\imgdir/#1.jpg}{%
    \begin{figure}[htbp]\centering
      \fcolorbox{gray!40}{white}{\includegraphics[width=0.96\textwidth]{\imgdir/#1.jpg}}
      \caption{#2}
    \end{figure}}{}}

\newtcolorbox{tip}[1]{enhanced, colback=goldlight, colframe=gold,
  boxrule=0.6pt, arc=3pt, left=8pt, right=8pt, top=5pt, bottom=5pt,
  fonttitle=\sffamily\bfseries, coltitle=black, colbacktitle=gold!35,
  title={#1}}
\newtcolorbox{note}[1]{enhanced, colback=violetlight, colframe=violet,
  boxrule=0.6pt, arc=3pt, left=8pt, right=8pt, top=5pt, bottom=5pt,
  fonttitle=\sffamily\bfseries, coltitle=white, colbacktitle=violet,
  title={#1}}

% Two-column tables that fill the text width.
\newenvironment{keytable}[2]{%
  \renewcommand{\arraystretch}{1.3}%
  \rowcolors{2}{white}{violetlight}%
  \tabularx{\textwidth}{@{}>{\raggedright\arraybackslash}p{0.33\textwidth}X@{}}
  \toprule \textsf{\textbf{#1}} & \textsf{\textbf{#2}} \\ \midrule}{%
  \bottomrule\endtabularx}

% Title page: \maketitlepage{subtitle}{version line}{blurb}
\newcommand{\maketitlepage}[3]{%
  \begin{titlepage}
    \centering
    \vspace*{3cm}
    \includegraphics[width=4.2cm]{../../assets/icon/app_icon.png}\par\vspace{1.2cm}
    {\sffamily\fontsize{40}{46}\selectfont\bfseries\color{violet} Illuminati\par}
    \vspace{0.5cm}
    {\sffamily\LARGE #1\par}
    \vspace{1.5cm}
    {\large\itshape #3\par}
    \vfill
    {\sffamily\color{gray} #2\par}
  \end{titlepage}}

\hypersetup{pdftitle={Illuminati — Manuale d’uso}, pdfauthor={Illuminati}, pdflang={it}}

\begin{document}

\maketitlepage{Manuale d’uso}{Versione 1.0 · Windows, macOS, Linux e Android}%
  {Trova, scarica e leggi libri mai coperti da diritto d’autore\\o i cui diritti sono scaduti.}

{\setcounter{tocdepth}{1}\tableofcontents}
\newpage

% ---------------------------------------------------------------------------
\section{Benvenuto}

Illuminati è un’app di lettura per Windows, macOS, Linux e Android. Funziona come le app per
ebook che forse conosci già, con una differenza: tutti i libri che offre sono gratuiti e legali da
scaricare, perché i loro diritti d’autore sono scaduti o non sono mai esistiti.

\begin{itemize}
  \item Oltre 75.000 libri da \textbf{Project Gutenberg}, più \textbf{Standard Ebooks},
        \textbf{Wikisource} e \textbf{Internet Archive}.
  \item Un lettore comodo con evidenziazioni, note, segnalibri, ricerca, dizionario e traduzione.
  \item La tua biblioteca personale, con scaffali e liste di lettura.
  \item Impostazioni di visualizzazione e accessibilità, compreso un carattere pensato per la
        dislessia.
  \item Nessun account, nessun abbonamento, nessuna pubblicità. Tutto resta sul tuo dispositivo.
\end{itemize}

L’interfaccia è disponibile in italiano, inglese, spagnolo, francese, tedesco e portoghese.

% ---------------------------------------------------------------------------
\section{Primi passi}

\subsection{Installazione}
\begin{itemize}
  \item \textbf{Windows:} copia l’intera cartella \texttt{Illuminati} sul computer e avvia
        \texttt{illuminati.exe}. Lascia gli altri file accanto; il programma ne ha bisogno.
  \item \textbf{Android:} installa il file \texttt{.apk}. Potrebbe essere necessario consentire
        l’installazione da questa origine nelle impostazioni del telefono.
  \item \textbf{macOS e Linux:} installa la versione preparata per il tuo sistema.
\end{itemize}

\subsection{La schermata iniziale}
Su un computer l’app si apre ingrandita. La schermata iniziale ha due parti:

\begin{enumerate}
  \item \ui{La tua biblioteca} (in alto): i libri che hai scaricato, a partire dall’ultimo letto,
        ciascuno con una barra di avanzamento. Il pulsante \ui{Biblioteca e scaffali} apre tutta
        la tua biblioteca.
  \item \ui{Scopri libri gratuiti} (in basso): una casella di ricerca, la lingua del catalogo, le
        categorie, una riga per ogni catalogo aggiuntivo e, in fondo, il catalogo di Project
        Gutenberg.
\end{enumerate}

\shot{home}{La schermata iniziale.}

I due pulsanti in alto a destra aprono \ui{Schermo e accessibilità} (sezione~\ref{sec:display}) e
\ui{Informazioni}, che spiega da dove provengono i libri e riporta le licenze dei caratteri.

% ---------------------------------------------------------------------------
\section{Trovare libri}

\subsection{Ricerca}
Scrivi in \ui{Cerca per titolo o autore}. La ricerca:
\begin{itemize}
  \item trova qualsiasi parte di un titolo o di un autore: \emph{miserab} trova
        \emph{I miserabili};
  \item ignora accenti e maiuscole: \emph{misérable} e \emph{MISERABLE} danno gli stessi
        risultati;
  \item interroga tutti i cataloghi insieme; ognuno mostra i risultati nella propria riga.
\end{itemize}
Svuota la casella con il pulsante \ui{×} per tornare ai libri popolari.

\subsection{Categorie e lingua}
Le etichette sotto la casella di ricerca (\ui{Narrativa}, \ui{Gialli}, \ui{Fantascienza},
\ui{Poesia}…) limitano tutti i cataloghi a un argomento. Il menu della lingua a destra mostra solo
i libri in una lingua, oppure \ui{Tutte le lingue}. I cataloghi senza risultati per la lingua o la
categoria scelta vengono nascosti.

\subsection{I cataloghi}
\begin{keytable}{Catalogo}{Cosa offre}
  Project Gutenberg & La più grande biblioteca gratuita, con oltre 75.000 libri ordinati per
                      popolarità. Molti hanno un’edizione illustrata. \\
  Standard Ebooks   & Classici revisionati con cura e ben impaginati, in inglese. \\
  Wikisource        & Testi trascritti da volontari, in molte lingue. L’ebook viene preparato al
                      momento del download; le opere molto lunghe (per esempio romanzi in più
                      volumi) potrebbero non essere disponibili. \\
  Internet Archive  & Scansioni di vecchie edizioni a stampa, pubblicate 95 o più anni fa. Il
                      testo è stato estratto automaticamente dalle scansioni e può contenere
                      errori. \\
\end{keytable}

\subsection{La pagina di un libro}
Tocca una copertina per aprire la pagina del libro. Mostra il riassunto, gli argomenti, la lingua,
il catalogo e una nota sul diritto d’autore. Da qui puoi:
\begin{itemize}
  \item \ui{Scarica gratis}: il libro viene salvato sul dispositivo, pronto da leggere offline.
        Ogni copertina del catalogo ha anche un piccolo pulsante di download.
  \item \ui{Leggi ora} o \ui{Continua a leggere}, una volta scaricato.
  \item \ui{Da leggere}, \ui{Letto} e \ui{Aggiungi a uno scaffale}, anche prima di scaricarlo
        (sezione~\ref{sec:library}).
  \item \ui{Scarica l’edizione illustrata}, per i libri di Gutenberg scaricati senza immagini.
\end{itemize}

\shot{details}{La pagina di un libro.}

% ---------------------------------------------------------------------------
\section{La tua biblioteca}\label{sec:library}

\ui{Biblioteca e scaffali}, nella schermata iniziale, apre tutti i libri che hai salvato: quelli
scaricati e quelli messi da parte per dopo senza scaricarli (con la copertina sbiadita e una
nuvoletta).

\shot{library}{La biblioteca, con le sue viste e i suoi scaffali.}

\subsection{Viste, filtro e ordinamento}
Le etichette in alto scelgono una vista, ognuna con il numero di libri: \ui{Tutti},
\ui{Scaricati}, \ui{Da leggere}, \ui{In lettura}, \ui{Letto} e poi ciascuno dei tuoi scaffali. La
casella sopra filtra per titolo o autore. Il pulsante di ordinamento (in alto a destra) ordina i
libri per \ui{Letti di recente}, \ui{Titolo}, \ui{Autore} (per cognome), \ui{Avanzamento} o
\ui{Aggiunti di recente}. L’app ricorda la tua scelta.

\subsection{Liste di lettura}
\begin{itemize}
  \item \ui{Da leggere} mette da parte un libro per dopo, senza scaricarlo. Quando apri il libro,
        esce da questa lista.
  \item \ui{Letto} viene impostato da solo quando arrivi alla fine di un libro (puoi annullarlo),
        oppure a mano.
\end{itemize}

\subsection{Scaffali}
Gli scaffali sono le tue raccolte personali, come \emph{Preferiti} o \emph{Filosofia}. Un libro
può stare su più scaffali.
\begin{itemize}
  \item Creane uno con \ui{Nuovo scaffale} nella biblioteca, oppure da
        \ui{Aggiungi a uno scaffale} su qualsiasi libro.
  \item Fai clic destro (o tieni premuto) su uno scaffale per rinominarlo o eliminarlo. Se elimini
        uno scaffale, i suoi libri restano nella biblioteca.
\end{itemize}

\subsection{Il menu di un libro}
Fai clic destro su un libro (o tienilo premuto sul telefono) per vedere \ui{Dettagli del libro},
\ui{Aggiungi a uno scaffale}, \ui{Da leggere}, \ui{Segna come letto} e due modi per rimuoverlo:
\begin{itemize}
  \item \ui{Rimuovi il download} libera spazio ma lascia il libro nella biblioteca, con note,
        scaffali e posizione di lettura.
  \item \ui{Rimuovi dalla biblioteca} elimina tutto, comprese note e segnalibri.
\end{itemize}

% ---------------------------------------------------------------------------
\section{Leggere}

Tocca un libro scaricato per aprirlo. Si apre dove l’avevi lasciato.

\shot{reader}{Il lettore con la sua barra degli strumenti.}

\subsection{Girare pagina}
\begin{keytable}{Azione}{Come}
  Pagina successiva  & Tocca il bordo destro · \kright{} · \kbd{Pag Giù} · \kbd{Spazio} \\
  Pagina precedente  & Tocca il bordo sinistro · \kleft{} · \kbd{Pag Su} \\
  Scorrere un po’    & \kup{} \kdown{} o la rotellina del mouse \\
  Mostrare la barra  & Tocca il centro della pagina \\
\end{keytable}

\subsection{La barra degli strumenti}
Tocca il centro della pagina per mostrarla. In alto: torna alla biblioteca, ricerca
(sezione~\ref{sec:search}), segnalibro, indice e impostazioni di visualizzazione. In basso:
capitolo precedente e successivo, un cursore per saltare tra i capitoli e il tuo avanzamento.

\subsection{Avanzamento e tempo rimanente}
La riga in basso indica quanti minuti mancano alla fine del capitolo e a che punto sei del libro.
Il tempo si basa sulla tua velocità di lettura, che l’app impara mentre leggi.

\subsection{Indice e immagini}
\ui{Indice} mostra i veri capitoli del libro, anche quando un unico file ne contiene diversi. Le
immagini compaiono nel testo; toccale per vederle a schermo intero e ingrandirle.

% ---------------------------------------------------------------------------
\section{Evidenziazioni, note e segnalibri}

\subsection{Selezionare il testo}
Sul telefono, tieni premuta una parola e trascina le maniglie. Sul computer, trascina il mouse sul
testo o fai doppio clic su una parola, poi fai \textbf{clic destro} sulla selezione. Compare un
menu con:
\begin{itemize}
  \item \ui{Evidenzia}: segna il testo con il tuo colore di evidenziazione.
  \item \ui{Nota}: un’evidenziazione con una tua nota, in uno di quattro colori.
  \item \ui{Definisci} e \ui{Traduci} (sezione~\ref{sec:lookup}).
  \item \ui{Copia}.
\end{itemize}
Le evidenziazioni si adattano a parole intere. Quelle con una nota sono sottolineate. Seleziona di
nuovo un passo evidenziato per modificarne la nota o rimuoverlo.

\subsection{Segnalibri}
Usa il pulsante segnalibro nella barra degli strumenti, oppure \kbd{Ctrl+B}. Un nastro rosso
indica che la pagina ha un segnalibro.

\subsection{Il pannello laterale}
Il pulsante indice apre un pannello con tre schede: \ui{Indice}, \ui{Segnalibri} e \ui{Note}.
Tocca una voce per andarci. \ui{Copia tutto in Markdown} copia tutte le evidenziazioni e le note,
pronte da incollare in un altro programma.

% ---------------------------------------------------------------------------
\section{Dizionario e traduzione}\label{sec:lookup}

Seleziona una parola (fino a tre) e scegli \ui{Definisci} per vederne i significati. Quelli nella
lingua del libro compaiono per primi. Provengono dal Wiktionary e sono scritti in inglese.

Seleziona un qualsiasi passo e scegli \ui{Traduci}. Scegli la lingua di destinazione dall’elenco;
l’app la ricorda.

\begin{note}{Serve internet}
Definizioni e traduzioni richiedono una connessione a internet. Il testo selezionato viene inviato
al Wiktionary o a MyMemory; nient’altro viene inviato.
\end{note}

% ---------------------------------------------------------------------------
\section{Cercare in un libro}\label{sec:search}

Premi \kbd{Ctrl+F} o il pulsante di ricerca nella barra degli strumenti e scrivi almeno due
lettere. La ricerca ignora maiuscole e accenti. Scegli un risultato per andarci. Una barra in alto
mostra \emph{3 di 27}: usa le sue frecce, o \kbd{F3}, per passare al risultato successivo.
\kbd{Esc} la chiude.

% ---------------------------------------------------------------------------
\section{Schermo e accessibilità}\label{sec:display}

Apri queste impostazioni con il pulsante delle impostazioni nella schermata iniziale, oppure con il
pulsante del testo nel lettore. Un’anteprima in alto mostra subito ogni modifica.

\shot{settings}{Schermo e accessibilità.}

\subsection{Lingua}
\ui{Lingua dell’app} cambia subito l’interfaccia in italiano, inglese, spagnolo, francese, tedesco
o portoghese. \ui{Come il dispositivo} segue la lingua del dispositivo.

\subsection{Testo}
\begin{itemize}
  \item Tema: \ui{Chiaro}, \ui{Seppia} o \ui{Scuro}.
  \item Carattere: quattro caratteri inclusi nell’app, così i libri hanno lo stesso aspetto su ogni
        dispositivo:
  \begin{itemize}
    \item \textbf{Literata}, progettato per leggere sullo schermo (predefinito);
    \item {\garamondfont EB Garamond}, un classico carattere da libro;
    \item {\sffamily Atkinson Hyperlegible}, creato perché le lettere simili si distinguano
          facilmente;
    \item {\dyslexicfont OpenDyslexic}, pensato per i lettori con dislessia.
  \end{itemize}
  \item \ui{Dimensione del testo}, \ui{Interlinea} e \ui{Spaziatura tra le lettere}.
  \item \ui{Giustifica il testo}: disattivalo per avere spazi uguali tra le parole, cosa che molti
        lettori con dislessia trovano più facile.
\end{itemize}

\subsection{Pagina}
\ui{Margini} e \ui{Larghezza delle righe}. Le righe più corte si seguono più facilmente.

\subsection{Schermo}
\ui{Luminosità} scurisce tutta l’app. \ui{Calore} riduce la luce blu, come una modalità notturna.
Entrambi si sommano alle impostazioni del dispositivo.

\ui{Ripristina predefiniti} riporta tutte queste impostazioni allo stato iniziale.

% ---------------------------------------------------------------------------
\section{Scorciatoie da tastiera}

\begin{keytable}{Tasti}{Azione}
  \kright{} \kbd{Pag Giù} \kbd{Spazio} & Pagina successiva \\
  \kleft{} \kbd{Pag Su}                & Pagina precedente \\
  \kup{} \kdown{}                      & Scorrere un po’ \\
  \kbd{Ctrl+F}                         & Cercare nel libro \\
  \kbd{F3}                             & Risultato successivo \\
  \kbd{Ctrl+B}                         & Aggiungere o togliere un segnalibro alla pagina \\
  \kbd{Esc}                            & Chiudere la ricerca o la barra, o tornare alla
                                         biblioteca \\
  \kbd{Maiusc} + frecce                & Estendere una selezione di testo \\
\end{keytable}

% ---------------------------------------------------------------------------
\section{Privacy e dati}

Illuminati non ha account né tracciamento. La biblioteca, le note e le impostazioni vengono
salvate solo sul tuo dispositivo (su Windows, in \texttt{\%APPDATA\%\textbackslash Illuminati}).

L’app si collega a internet solo per:
\begin{itemize}
  \item sfogliare i cataloghi, cercarvi libri e scaricare libri e copertine;
  \item definire o tradurre il testo selezionato, che viene inviato al Wiktionary o a MyMemory.
\end{itemize}

% ---------------------------------------------------------------------------
\section{Sul diritto d’autore}

Tutti i libri sono di pubblico dominio negli Stati Uniti. Alcuni testi di Wikisource hanno invece
licenze libere (come CC BY-SA), che consentono di leggerli e condividerli liberamente. Da Internet
Archive vengono mostrati solo libri pubblicati 95 o più anni fa.

\begin{tip}{Controlla le leggi del tuo paese}
Il diritto d’autore cambia da paese a paese. Un libro libero negli Stati Uniti potrebbe essere
ancora protetto dove vivi, quindi controlla le leggi del tuo paese prima di scaricare.
\end{tip}

% ---------------------------------------------------------------------------
\section{Risoluzione dei problemi}

\begin{keytable}{Messaggio}{Cosa fare}
  Impossibile raggiungere un catalogo & Il servizio potrebbe essere sovraccarico o la
                             connessione assente. Usa \ui{Riprova} più tardi; gli altri cataloghi
                             continuano a funzionare. \\
  Questo titolo non ha un’edizione ebook & Probabilmente è un audiolibro. Cerca un’altra
                             edizione. \\
  Wikisource non è riuscita a preparare questo libro & L’opera è troppo lunga per il servizio di
                             Wikisource. Prova uno dei suoi volumi o un altro catalogo. \\
  Raggiunto il limite giornaliero di traduzioni gratuite & MyMemory consente un numero limitato di
                             traduzioni gratuite al giorno. Riprova domani. \\
  Il download ha impiegato troppo tempo & Controlla la connessione e riprova. \\
\end{keytable}

% ---------------------------------------------------------------------------
\section{Crediti}

Libri: Project Gutenberg, Standard Ebooks, Wikisource (tramite WS Export) e Internet Archive.
Definizioni: Wiktionary (CC BY-SA). Traduzioni: MyMemory.

Caratteri con licenza SIL Open Font License: Literata, EB Garamond, Atkinson Hyperlegible (Braille
Institute) e OpenDyslexic (Abbie Gonzalez). Le loro licenze sono in \ui{Informazioni} →
\ui{Licenze} nell’app.

\end{document}
